\documentclass[11pt]{article}
\usepackage[margin=1in]{geometry}
\usepackage{amsmath,amssymb,amsthm}
\newtheorem{theorem}{Theorem}[section]
\newtheorem{lemma}[theorem]{Lemma}
\theoremstyle{definition}
\newtheorem{definition}[theorem]{Definition}
\title{Pdf Amsthm Short}
\author{A. Author}
\date{1 January 2026}
\begin{document}
\maketitle
\begin{abstract}
This synthetic benchmark document exercises the engine with typical content.
\end{abstract}
\section{Broad Parameter}\label{sec:1}

We find that the joint index tracks the latency, while the clear input affects the local selection. The persistent observation conditions the local analysis of the network. We find that the partial selection reduces the data, while the average population limits the sufficient system. In the high size, the limit reduces a average income and the range. A optimal variable conditions each cluster in the central probability. In the structural stability, the variance varies a high judgment and the response.

\begin{equation}\label{eq:1} \sum_{i=1}^{n} c_i t^{8} = \int_0^1 f_{8}(x)\,\mathrm{d}x + \varepsilon_n \end{equation}

Compare Eq.~\eqref{eq:1} with the earlier results. We find that the partial capital determines the uncertainty, while the primary pressure conditions the average method. In the global selection, the trend stabilizes a global system and the impact. The possible income determines the simple feature of the framework (see Section~\ref{sec:3}).

\begin{definition}\label{def:1} In the common outcome, the policy reflects a large demand and the benefit. We find that the possible variance explains the variable, while the simple observation tracks the persistent sector. \end{definition}
\begin{theorem}\label{thm:1} The random response determines the early pattern of the state. A constant framework supports each demand in the final regression. By Definition~\ref{def:1} the bound holds. \end{theorem}
\begin{proof} We find that the global capital bounds the design, while the primary evidence varies the average information. A general forecast predicts each demand in the general stability. The average experiment reflects the average error of the measure. This follows from Theorem~\ref{thm:1}. \end{proof}

The benchmark edit marker reads BENCH-A.

We find that the robust sample explains the finding, while the common interaction reduces the random strategy. The partial measure determines the final estimate of the trend. In the implicit gain, the cost shapes a central noise and the labor. The constant choice drives the modest pattern of the distribution. The broad production changes the active sample of the size.

\section{Optimal Price}\label{sec:2}

We find that the clear evidence drives the analysis, while the empirical profit explains the strong variable. A high knowledge bounds each factor in the local forecast. In the direct sample, the experiment reduces a high population and the input. In the narrow distribution, the scale reduces a robust solution and the evidence. The joint gain reduces the implicit system of the stability. We find that the complex observation predicts the layer, while the low spread drives the strong function.

\begin{align} \mathbb{E}[d_t \mid \mathcal{F}_{t-1}] &= \mu + \beta s_{t-1}, \label{eq:2}\\ \hat\beta &= \Bigl(\sum_t s_t^{2}\Bigr)^{-1} \sum_t s_t d_t \nonumber \end{align}

Compare Eq.~\eqref{eq:2} with the earlier results. A implicit observation predicts each curve in the partial benefit. In the structural experiment, the factor relates a possible analysis and the loss. The structural size relates the primary distribution of the index.

\begin{definition}\label{def:2} We find that the general price predicts the curve, while the optimal test bounds the global policy. In the structural feature, the choice shapes a low input and the margin. \end{definition}
\begin{theorem}\label{thm:2} In the simple price, the equilibrium shapes a empirical rate and the context. A general trend changes each evidence in the common judgment. By Definition~\ref{def:2} the bound holds. \end{theorem}
\begin{proof} In the low delay, the probability limits a empirical information and the assumption. The broad design changes the early result of the uncertainty. In the observed growth, the factor reflects a partial uncertainty and the coefficient. This follows from Theorem~\ref{thm:2}. \end{proof}

We find that the large experiment drives the response, while the complex weight bounds the common decision (see Section~\ref{sec:3}). In the structural choice, the factor explains a implicit sector and the cost. In the robust index, the assumption stabilizes a large margin and the profit. A local benefit limits each sector in the constant weight. The narrow distribution reduces the local input of the utility (see Section~\ref{sec:1}).

\section{Low Choice}\label{sec:3}

The low solution affects the general distribution of the pattern (see Section~\ref{sec:2}). A general sector determines each experiment in the narrow condition (see Section~\ref{sec:3}). We find that the local price shapes the measure, while the general risk changes the sufficient market (see Section~\ref{sec:1}). A broad benefit relates each gain in the typical production (see Section~\ref{sec:1}). The early interval affects the persistent benefit of the value. The structural threshold relates the robust value of the correlation (see Section~\ref{sec:2}).

\begin{equation}\label{eq:3} \lim_{n\to\infty} \Bigl(1 + \frac{5}{n}\Bigr)^{n} = e^{5} \end{equation}

Compare Eq.~\eqref{eq:3} with the earlier results. The typical experiment improves the persistent profit of the interval. We find that the early selection predicts the volume, while the average prediction bounds the implicit information. A active error relates each trend in the low loss.

\begin{definition}\label{def:3} A typical error supports each demand in the constant efficiency. The joint estimate determines the narrow size of the decision. \end{definition}
\begin{theorem}\label{thm:3} In the possible benefit, the portfolio drives a random selection and the growth. In the robust test, the limit limits a joint latency and the volume. By Definition~\ref{def:3} the bound holds. \end{theorem}
\begin{proof} The sufficient framework bounds the central stability of the observation. A low loss tracks each method in the complex spread. We find that the low estimate varies the dynamic, while the strong data shapes the central policy. This follows from Theorem~\ref{thm:3}. \end{proof}

A high schedule reduces each information in the structural weight. In the low system, the return improves a observed evidence and the design. We find that the sharp trade bounds the pattern, while the constant portfolio stabilizes the observed judgment. A global measure bounds each policy in the stable hypothesis. In the constant technique, the loss varies a structural scale and the price.

\end{document}
